% GENERATED FILE. Edit canonical lessons, then run npm run generate:books.
% canonical-source-hash: fnv1a64:55da8d7886e63eee
% canonical-lessons: ES-C774-sotano, ES-C774-desvan, ES-C774-cajon, ES-C774-espejo, ES-C774-persiana

\chapter{La casa de mis abuelos}
\label{ch:es-la-casa-de-mis-abuelos}
\hlchaptermodality{774}

\begin{chapteropening}
\textbf{By the end of this chapter:} \emph{I can describe an old house as I remember it: the cellar, the attic, a drawer, a mirror and the blinds.}

Complete the last lesson of chapter 774: I can describe an old house as I remember it: the cellar, the attic, a drawer, a mirror and the blinds.
\end{chapteropening}

\section[el sótano]{el sótano — the basement, the cellar}

\label{lesson:ES-C774-sotano}



\begin{quote}
Before the new one: say the Spanish for the link, then the Spanish for to top up.
\end{quote}

\subsection*{You'll want to know: el sótano}
\textbf{el sótano} — \textquotedblleft{}the basement, the cellar\textquotedblright{}.

\begin{quote}
El sótano siempre estaba frío y oscuro.
\end{quote}

\begin{quote}
\emph{The basement was always cold and dark.}
\end{quote}

\begin{cousinweb}
\emph{sótano} is usually traced to spoken Latin \emph{*subtulum} 'a place below', built on \emph{subtus} 'beneath' (Lewis \& Short).
\end{cousinweb}

\begin{grammarlens}[title={What you've built}]
One more word for telling, explaining and describing at B1.
\end{grammarlens}

\subsection*{Your turn}
\begin{itemize}
\raggedright
  \item \emph{Say it:} \emph{el sótano}
  \item \emph{Say it:} \emph{el sótano}, once more
  \item \emph{Say it:} a sentence of your own with \emph{el sótano}
  \item \emph{Recall:} say the Spanish for the link, then the Spanish for to top up, then say \emph{el sótano} again
\end{itemize}

\begin{tcolorbox}[breakable,colback=teal!4,colframe=teal!35!black,title={Before you move on}]
What does el sótano mean? (\textquotedblleft{}The basement, the cellar\textquotedblright{}.) Say it once more.
\end{tcolorbox}
\section[el desván]{el desván — the attic, the loft}

\label{lesson:ES-C774-desvan}



\begin{quote}
Before the new one: say the Spanish for to top up, then the Spanish for the basement.
\end{quote}

\subsection*{You'll want to know: el desván}
\textbf{el desván} — \textquotedblleft{}the attic, the loft\textquotedblright{}.

\begin{quote}
En el desván había cajas de fotos viejas.
\end{quote}

\begin{quote}
\emph{In the attic there were boxes of old photos.}
\end{quote}

In Spain \emph{el ático} is usually a top-floor flat, not a loft.

\begin{cousinweb}
\emph{desván} is usually traced to \emph{des-} and \emph{vano} 'an empty space', from Latin \emph{vānus} 'empty, void' (Lewis \& Short); the details are not settled.
\end{cousinweb}

\begin{grammarlens}[title={What you've built}]
One more word for telling, explaining and describing at B1.
\end{grammarlens}

\subsection*{Your turn}
\begin{itemize}
\raggedright
  \item \emph{Say it:} \emph{el desván}
  \item \emph{Say it:} \emph{el desván}, once more
  \item \emph{Say it:} a sentence of your own with \emph{el desván}
  \item \emph{Recall:} say the Spanish for to top up, then the Spanish for the basement, then say \emph{el desván} again
\end{itemize}

\begin{tcolorbox}[breakable,colback=teal!4,colframe=teal!35!black,title={Before you move on}]
What does el desván mean? (\textquotedblleft{}The attic, the loft\textquotedblright{}.) Say it once more.
\end{tcolorbox}
\section[el cajón]{el cajón — the drawer}

\label{lesson:ES-C774-cajon}



\begin{quote}
Before the new one: say the Spanish for the basement, then the Spanish for the attic.
\end{quote}

\subsection*{You'll want to know: el cajón}
\textbf{el cajón} — \textquotedblleft{}the drawer\textquotedblright{}.

\begin{quote}
Mi abuela guardaba las cartas en un cajón.
\end{quote}

\begin{quote}
\emph{My grandmother kept the letters in a drawer.}
\end{quote}

\emph{De cajón} — obvious, of course.

\begin{cousinweb}
\emph{cajón} is built on \emph{caja} 'box' with the ending \emph{-ón}. \emph{Caja} comes from Latin \emph{capsa} 'a box, a case for books' (Lewis \& Short), most likely through Catalan \emph{caixa}.
\end{cousinweb}

\begin{grammarlens}[title={What you've built}]
One more word for telling, explaining and describing at B1.
\end{grammarlens}

\subsection*{Your turn}
\begin{itemize}
\raggedright
  \item \emph{Say it:} \emph{el cajón}
  \item \emph{Say it:} \emph{el cajón}, once more
  \item \emph{Say it:} a sentence of your own with \emph{el cajón}
  \item \emph{Recall:} say the Spanish for the basement, then the Spanish for the attic, then say \emph{el cajón} again
\end{itemize}

\begin{tcolorbox}[breakable,colback=teal!4,colframe=teal!35!black,title={Before you move on}]
What does el cajón mean? (\textquotedblleft{}The drawer\textquotedblright{}.) Say it once more.
\end{tcolorbox}
\section[el espejo]{el espejo — the mirror}

\label{lesson:ES-C774-espejo}



\begin{quote}
Before the new one: say the Spanish for the attic, then the Spanish for the drawer.
\end{quote}

\subsection*{You'll want to know: el espejo}
\textbf{el espejo} — \textquotedblleft{}the mirror\textquotedblright{}.

\begin{quote}
En el pasillo había un espejo enorme.
\end{quote}

\begin{quote}
\emph{In the hall there was a huge mirror.}
\end{quote}

\emph{Mirarse en el espejo} — to see yourself in the mirror.

\begin{cousinweb}
\emph{espejo} is inherited from Latin \emph{speculum} 'a mirror' (Lewis \& Short), built on \emph{specere} 'to look' (EtymDB).
\end{cousinweb}

\begin{grammarlens}[title={What you've built}]
One more word for telling, explaining and describing at B1.
\end{grammarlens}

\subsection*{Your turn}
\begin{itemize}
\raggedright
  \item \emph{Say it:} \emph{el espejo}
  \item \emph{Say it:} \emph{el espejo}, once more
  \item \emph{Say it:} a sentence of your own with \emph{el espejo}
  \item \emph{Recall:} say the Spanish for the attic, then the Spanish for the drawer, then say \emph{el espejo} again
\end{itemize}

\begin{tcolorbox}[breakable,colback=teal!4,colframe=teal!35!black,title={Before you move on}]
What does el espejo mean? (\textquotedblleft{}The mirror\textquotedblright{}.) Say it once more.
\end{tcolorbox}
\section[la persiana]{la persiana — the blind, the shutter}

\label{lesson:ES-C774-persiana}



\begin{quote}
Before the new one: say the Spanish for the drawer, then the Spanish for the mirror.
\end{quote}

\subsection*{You'll want to know: la persiana}
\textbf{la persiana} — \textquotedblleft{}the blind, the shutter\textquotedblright{}.

\begin{quote}
Mi abuela siempre bajaba la persiana por la tarde.
\end{quote}

\begin{quote}
\emph{My grandmother always lowered the blind in the afternoon.}
\end{quote}

\emph{Subir} and \emph{bajar la persiana} — to raise and lower the blind.

\begin{cousinweb}
\emph{persiana} comes from French \emph{persienne} (EtymDB), a word built on the name of Persia; the slatted blind was thought to come from there.
\end{cousinweb}

\begin{grammarlens}[title={What you've built}]
One more word for telling, explaining and describing at B1.
\end{grammarlens}

\subsection*{Your turn}
\begin{itemize}
\raggedright
  \item \emph{Say it:} \emph{la persiana}
  \item \emph{Say it:} \emph{la persiana}, once more
  \item \emph{Say it:} a sentence of your own with \emph{la persiana}
  \item \emph{Recall:} say the Spanish for the drawer, then the Spanish for the mirror, then say \emph{la persiana} again
\end{itemize}

\begin{tcolorbox}[breakable,colback=teal!4,colframe=teal!35!black,title={Before you move on}]
What does la persiana mean? (\textquotedblleft{}The blind, the shutter\textquotedblright{}.) Say it once more.
\end{tcolorbox}
